\section{模型评价与推广}

\subsection{模型优点}

\begin{enumerate}
  \item 覆盖宽度与重叠率全程使用“水平投影”口径，可与水平测线间距直接比较，避免了斜坡情形下坡面长度与水平距离混用的问题；
  \item 问题一、二均为确定性解析几何模型，参数可解释、结果可复算，并在 $\alpha=0$ 或 $\beta=0^\circ,180^\circ$ 等特殊情形下退化为平底公式，模型正确性有明确检验；
  \item 问题三的贪心递推针对单一平面坡面可证明在当前方向下测线数不可再减少，结果具有可解释的最优性；
  \item 问题四采用分带策略并保留网格覆盖仿真，能够处理真实地形下的深浅差异，方案与题目三项指标直接对应。
\end{enumerate}

\subsection{模型不足}

\begin{enumerate}
  \item 问题一至问题三假设海底为理想平面斜坡、声波沿直线传播，未考虑真实声速剖面、折射和复杂地形；
  \item 问题四仅统计测量线段长度，未计入不同线段之间的转场、调头和航行连接成本，可能低估实际作业代价；
  \item 分带边界处各带独立布线的覆盖近似会引入少量重叠误差，局部最大相邻重叠率仍达 $49.01\%$；
  \item 问题四未进行随机扰动或全局优化，所得为当前分带与递推规则下的较优方案，而非严格全局最优。
\end{enumerate}

\subsection{模型推广}

模型可推广到一般起伏海底：保留“水平投影覆盖宽度 + 相邻重叠率”的评价框架，将平面坡度替换为局部梯度估计即可用于非规则地形。若进一步引入声速分层，可在射线追踪中修正声线曲率；若考虑船舶转弯与多船协同，可把分带递推扩展为带转场代价的路径优化问题。
